% =========================================================================
% Chapter: Literature Review
% Dissertation: Computer Vision for Urban Waste and Infrastructure
%               Monitoring in Malta (University of Malta, 2026)
% Written against repository commit 1b2c918a (inspected 11 July 2026);
% revised against commit 2d70be27 (inspected 16 July 2026) to add
% LingBot-Vision coverage.
% Literature search current through 16 July 2026.
% Citation style: natbib author-year (\citet / \citep).
% =========================================================================

\chapter{Literature Review}
\label{ch:literature-review}

This chapter critically reviews the published work most relevant to the dissertation. It proceeds from the application domain (computer vision for municipal monitoring, and the waste and traffic-sign datasets that support it) through the three technical families compared in this work: supervised object detection, prompt-based and open-vocabulary localisation, and self-supervised representation learning with its adaptation strategies. Throughout, the emphasis is on comparing approaches, weighing the evidence behind published claims, and identifying what the literature does not yet establish; the chapter closes with a synthesis of the research gaps that this dissertation addresses.

\section{Computer Vision for Municipal Monitoring}
\label{sec:lr-municipal}

Applied work on automated urban inspection has grown steadily with the availability of street-level imagery and affordable capture hardware. Early systems demonstrated feasibility on narrow slices of the problem: \citet{mittal2016spotgarbage} trained a convolutional network to flag garbage in crowdsourced smartphone photographs, and \citet{rad2017computer} mounted a camera on a street-sweeping vehicle to localise cigarette butts, leaves and packaging on Swiss streets with an early one-stage detector. On the infrastructure side, \citet{balali2015detection} demonstrated city-scale traffic-sign inventory from Google Street View imagery, and \citet{tabernik2020deep} showed that a Mask R-CNN-based pipeline could support inventory management across 200 sign categories on Slovenian roads. More recently, \citet{merolla2025improving} deployed YOLO-based detection of damaged signs and road-surface defects in an Italian regional pilot, explicitly framing the system as decision support for municipal maintenance prioritisation.

Read together, these studies establish that street-level detection for municipal purposes is technically feasible, but they also expose recurring weaknesses. Most report results on imagery from a single city or country, leaving geographical transfer untested---a concern given documented dataset bias effects \citep{torralba2011unbiased} and evidence that recognition systems degrade on visual environments unlike their training distribution \citep{devries2019does}. Small, distant objects remain the dominant error mode, consistent with the broader small-object detection literature \citep{cheng2023towards}. Class imbalance, occlusion by parked vehicles and illumination variation are frequently acknowledged but rarely analysed systematically. Deployment cost, edge feasibility and privacy handling receive limited treatment: few papers report latency or memory on realistic hardware, and fewer engage with data-protection constraints on street imagery, despite faces and number plates being unavoidable in this domain \citep{gdpr2016}. No published study identified in this search evaluates modern detection models on Maltese street imagery in either the waste or the signage domain.

\section{Waste Datasets and Waste Recognition}
\label{sec:lr-waste}

The waste-recognition literature is best organised by capture setting, because the setting determines how transferable results are to kerbside monitoring. At the controlled end, TrashNet contains single items photographed on a clean background for six-way material classification \citep{yang2016trashnet}; it remains widely used despite being a non-peer-reviewed student dataset, and results obtained on it say little about street scenes. Facility-oriented datasets such as ZeroWaste capture cluttered conveyor streams inside recycling plants and motivate deformable-object segmentation \citep{bashkirova2022zerowaste}, but industrial sorting imagery---top-down, controlled lighting, no occluding vehicles---differs fundamentally from uncontrolled urban scenes, and this review deliberately does not treat the extensive sorting literature \citep[surveyed by][]{lu2022computer} as directly relevant evidence for kerbside detection. Closer to the target setting, TACO annotates litter in diverse outdoor contexts with segmentation masks and a material-based taxonomy \citep{proenca2020taco}, and UAVVaste targets litter in aerial imagery \citep{kraft2021uavvaste}. Systematic reviews of artificial intelligence in solid-waste management confirm that classification of loose items and facility-level sorting dominate the literature, while street-level applications remain comparatively rare \citep{abdallah2020artificial}.

Two structural gaps emerge. First, existing datasets almost universally label \emph{materials} (plastic, metal, paper) or generic \emph{litter}, whereas municipal operations reason about \emph{collection streams}: a black residual bag, a white organic bag, a grey recycling bag. Recognising a sealed bag's stream from its presentation is a different problem from recognising exposed materials, and the reviewed literature does not provide a dataset or benchmark for it. Second, waste presentation is policy-dependent and therefore geographically specific: bag colours, container types and kerbside conventions differ across countries and even municipalities, so a model trained on one jurisdiction's presentation conventions has no guarantee of transferring to another's. Both gaps motivate a Malta-specific kerbside dataset organised around the national separation streams.

\section{Traffic-Sign Datasets and Recognition}
\label{sec:lr-signs}

Traffic-sign recognition is among the oldest benchmark-driven areas of road-scene understanding \citep{mogelmose2012vision}. The German Traffic Sign Recognition Benchmark (GTSRB) established cropped-sign classification, with early deep networks exceeding human accuracy \citep{stallkamp2012man}, and its companion detection benchmark GTSDB moved to full-scene localisation \citep{houben2013detection}. Subsequent datasets scaled scene realism and coverage: the LISA dataset introduced US signage and video \citep{mogelmose2012vision}, Tsinghua-Tencent 100K (TT100K) emphasised small signs in large panoramas \citep{zhu2016traffic}, the Mapillary Traffic Sign Dataset extended coverage to over 300 sign classes across six continents using crowdsourced street imagery \citep{ertler2020mapillary}, and CURE-TSD systematically injected challenging conditions such as rain, haze and blur to expose robustness failures \citep{temel2020traffic}. Task variants beyond detection and classification include sign inventory from mobile mapping \citep{balali2015detection}, large-category recognition for maintenance \citep{tabernik2020deep}, and, more recently, explicit damage and condition assessment \citep{merolla2025improving}.

Against this dissertation's requirements, three limitations of the established datasets stand out. First, taxonomy transfer: Maltese signage mixes European conventions with locally specific categories---bilingual street nameplates, tourist signage, auxiliary plates and convex blind-spot mirrors---that are either absent from, or inconsistently represented in, international datasets; models trained on GTSRB- or Mapillary-style taxonomies cannot emit these categories at all without retraining. Second, condition annotation is scarce. GTSRB, GTSDB, TT100K and Mapillary annotate identity but not physical state; CURE-TSD varies \emph{imaging} conditions rather than sign \emph{condition}; and the emerging condition-assessment literature typically annotates a binary damaged/undamaged distinction on small bespoke datasets \citep{merolla2025improving}. No established public dataset identified in this search jointly provides condition, mounting type, viewing angle and shape annotations per sign, which is precisely the combination required for maintenance-oriented monitoring. Third, the viewing-angle dimension is essentially unexplored: benchmark datasets overwhelmingly annotate front-facing signs, whereas inventory imagery captured from a moving vehicle unavoidably contains oblique and rear views that still carry asset information.

\section{Evolution of Supervised Object Detection}
\label{sec:lr-detection}

Modern detection descends from R-CNN, which classified externally proposed regions with a convolutional network \citep{girshick2014rich}; Fast R-CNN moved feature extraction before proposal pooling \citep{girshick2015fast}, and Faster R-CNN internalised proposal generation in a region proposal network, defining the canonical \emph{two-stage} detector \citep{ren2015faster}. One-stage detectors removed the proposal stage entirely: YOLO framed detection as a single dense regression \citep{redmon2016you}, SSD added multi-scale prediction \citep{liu2016ssd}, and RetinaNet's focal loss addressed the extreme foreground--background imbalance that had kept dense detectors behind two-stage accuracy \citep{lin2017focal}, aided by feature pyramid necks \citep{lin2017feature}. The empirical trade-off that emerged---two-stage detectors offering higher accuracy on small and crowded objects, one-stage detectors offering lower latency and simpler deployment---narrowed over successive YOLO generations \citep{redmon2018yolov3,bochkovskiy2020yolov4,wang2023yolov7} until one-stage models became the default for real-time applications.

DETR reformulated detection as set prediction with a Transformer encoder--decoder and bipartite matching, eliminating hand-crafted anchors and NMS \citep{carion2020end}, at the cost of slow convergence and weak small-object performance; Deformable DETR's sparse multi-scale attention substantially repaired both \citep{zhu2021deformable}, and DINO-style improvements to query initialisation and denoising training brought DETR-family models to state-of-the-art accuracy \citep{zhang2023dino}. (Throughout this dissertation, ``DINO'' as a \emph{detector} \citep{zhang2023dino} must be distinguished from ``DINO'' as a \emph{self-supervised representation method} \citep{caron2021emerging}; the similar names denote unrelated systems.) RT-DETR then demonstrated that a carefully designed hybrid encoder could bring DETR-family accuracy into the real-time regime, reporting higher COCO AP than similarly sized YOLO models under the authors' evaluation configuration \citep{zhao2024detrs}, and LW-DETR showed that a plain ViT backbone with a shallow decoder could compete with YOLO at real-time latencies \citep{chen2024lwdetr}. The relevant axis for this dissertation is thus no longer ``CNN versus Transformer'' as a binary, but a spectrum of design choices---dense prediction with duplicate suppression versus sparse set prediction, convolutional versus attention-based feature extraction, and from-scratch versus pretrained initialisation---each with measurable consequences for small objects, crowded scenes, training cost and edge deployment.

\section{The YOLO Generations under Study}
\label{sec:lr-yolo}

The dissertation's supervised track uses three recent YOLO generations, all released by or benchmarked against Ultralytics tooling, and their published characteristics deserve individual scrutiny.

\textbf{YOLO11} (Ultralytics, September 2024) is a convolution-dominated anchor-free detector refining the C3k2 backbone blocks, SPPF pooling and decoupled head lineage of its predecessors, released in five scales (n--x) with multi-task support \citep{ultralytics2024yolo11}. No peer-reviewed paper accompanies it; the most-cited description is a third-party architectural overview \citep{khanam2024yolov11}, and its reported COCO improvements over YOLOv8 at reduced parameter counts derive from the vendor's own benchmarks, which should be treated as indicative rather than independently validated.

\textbf{YOLO12} \citep{tian2025yolov12} is, by contrast, an academically authored generation with an explicit architectural thesis: that attention can be made fast enough for real-time detection. Its area-attention mechanism restricts self-attention to large image regions, preserving receptive field while cutting cost, supported by residual efficient-layer aggregation blocks and flash-attention kernels. The authors report YOLO12-N at 40.6\% COCO mAP with 1.64\,ms T4 latency, exceeding YOLO11-N by 1.2 points at comparable speed, and favourable comparisons against RT-DETR variants. Two caveats matter for this dissertation: the reported gains are modest at matched scales, and the attention-centric design's dependence on modern GPU kernels means its speed advantages may not survive translation to CPU-only or embedded municipal hardware.

\textbf{YOLO26} (Ultralytics, released January 2026) is the most consequential and least independently validated of the three. Its technical report describes a natively end-to-end, NMS-free design with a dual-head scheme that also supports conventional dense prediction, removal of distribution focal loss in favour of a lighter unconstrained regression head, progressive loss balancing and small-target-aware label assignment, and a hybrid MuSGD optimiser imported from language-model training; the vendor reports 40.9--57.5 COCO mAP across five scales at 1.7--11.8\,ms TensorRT latency and up to 43\% faster CPU inference than YOLO11, with edge deployment as an explicit design goal \citep{jocher2026yolo26,ultralytics2026yolo26docs}. These claims currently rest on the vendor's own preprint and documentation; the sole independent treatment identified is an academic review that analyses the architecture and deployment pathways but does not re-measure the headline benchmarks \citep{sapkota2025yolo26}. Its NMS-free formulation follows the one-to-one assignment direction opened by YOLOv10 \citep{wang2024yolov10}, which matters practically because it removes a latency-variable post-processing stage and a tunable IoU threshold from deployment. For all three generations, COCO performance is an imperfect proxy for MDWD and MTSD performance: COCO's object statistics, class balance and scene composition differ substantially from cluttered Maltese street imagery, and the transfer of vendor-reported margins to these domains is an open empirical question that the dissertation's experiments address directly.

\section{Non-Maximum Suppression and End-to-End Detection}
\label{sec:lr-nms}

Dense detectors necessarily produce multiple overlapping predictions per object, because many spatial locations independently score the same instance; classical NMS resolves this by greedily keeping the highest-scoring box and suppressing neighbours above an IoU threshold \citep{neubeck2006efficient}. The procedure is simple but consequential: it adds a hyperparameter whose optimal value is dataset-dependent, its cost grows with prediction count, and in crowded scenes it suppresses genuinely distinct adjacent objects---the classic failure being pedestrians in a crowd \citep{bodla2017soft}. Soft-NMS decays rather than removes neighbouring scores \citep{bodla2017soft}, and learned suppression networks re-score detections jointly \citep{hosang2017learning}, but both retain a separate post-processing stage. Set-prediction training removes the problem at its source: bipartite matching penalises duplicate predictions during training, so the network itself learns one-to-one assignment \citep{carion2020end}, an idea transplanted into the YOLO family through YOLOv10's dual-assignment training \citep{wang2024yolov10} and adopted natively by YOLO26 \citep{jocher2026yolo26}. This distinction is not academic for municipal imagery: kerbside waste bags are routinely placed touching one another in piles, and traffic signs cluster on shared poles, so the boundary between ``duplicate prediction'' and ``adjacent instance'' is exactly where NMS-based and end-to-end detectors can be expected to behave differently. The literature, however, evaluates NMS-free designs almost exclusively on COCO-style benchmarks; evidence on grouped-bag or clustered-sign configurations is not available, leaving a question this dissertation's per-class and failure-case analyses can inform.

\section{Real-Time Detection Transformers and RF-DETR}
\label{sec:lr-rfdetr}

RF-DETR \citep{robinson2025rfdetr} is the transformer-family detector included in the dissertation's supervised track. Architecturally it combines a DINOv2-with-registers ViT backbone with windowed attention, a multi-scale projector, and a Deformable-DETR-style decoder, drawing directly on LW-DETR's lightweight-decoder lineage \citep{chen2024lwdetr,zhu2021deformable,oquab2024dinov2}; training-wise, its distinctive contribution is weight-sharing neural architecture search that fine-tunes a pretrained base network once and then evaluates thousands of accuracy--latency configurations without retraining. The authors report that RF-DETR is the first real-time detector to exceed 60~mAP on COCO and that it outperforms LW-DETR and YOLO11 on both COCO and the RF100-VL multi-domain benchmark; the work is peer reviewed (ICLR 2026), and its use of a self-supervised foundation backbone is explicitly credited with strong adaptation to small, domain-specific datasets. Several qualifications temper these claims for present purposes. First, scale: published RF-DETR sizes span roughly 29M--128M parameters, so a comparison against nano- or small-scale YOLO models is a comparison across capacity classes, and any cross-family conclusion must name the scales involved. Second, the pretrained-backbone advantage confounds architecture with initialisation: RF-DETR inherits DINOv2's web-scale pretraining, whereas YOLO baselines are typically COCO-pretrained only, so ``transformer beats CNN'' readings of its results are not supported without controlling for pretraining. Third, early third-party comparisons in agricultural imagery found RF-DETR stronger on single-class ambiguity but YOLOv12 stronger under some multi-class regimes, suggesting domain-dependent trade-offs rather than dominance. Its behaviour on small, cluttered Maltese municipal imagery is unestablished.

\section{Transfer Learning versus Training from Scratch}
\label{sec:lr-transfer}

The default practice of initialising detectors and classifiers from large-corpus pretraining rests on well-replicated evidence: features learned on ImageNet-scale data transfer broadly \citep{yosinski2014transferable}, better pretrained classifiers tend to transfer better \citep{kornblith2019better}, and pretraining buys large sample-efficiency gains on small downstream datasets. The evidence is, however, more conditional than the practice suggests. \citet{he2019rethinking} showed that on COCO-scale data, training from random initialisation matches fine-tuned accuracy given sufficient schedule length, with pretraining mainly accelerating convergence; conversely, for datasets of a few thousand images---the regime of both MDWD's unique source images and MTSD---random initialisation is rarely competitive. Domain shift moderates the benefit: pretraining distributions dominated by web photography transfer imperfectly to specialised viewpoints and rare categories, and source-domain biases can persist through fine-tuning \citep{torralba2011unbiased,devries2019does}. For this dissertation the practical questions are therefore not \emph{whether} to use pretrained weights but \emph{which} pretraining (supervised COCO detection weights versus self-supervised web-scale features) and \emph{how much} adaptation each model family needs---questions on which the detector literature and the representation-learning literature reviewed below give partially conflicting answers.

\section{Edge-Oriented Detection}
\label{sec:lr-edge}

A parallel literature addresses making detectors deployable on constrained hardware: efficiency-first backbone design \citep{howard2017mobilenets,sandler2018mobilenetv2}, pruning and weight compression \citep{han2016deep}, integer quantisation with quantisation-aware training \citep{jacob2018quantization}, and knowledge distillation from larger teachers \citep{hinton2015distilling}, typically combined with compilation stacks such as TensorRT and reduced input resolutions. Surveys of edge inference consistently note that published ``real-time'' claims are only interpretable jointly with hardware, batch size, input resolution, numeric precision and whether pre- and post-processing are included \citep{chen2019deep,murshed2021machine}---conditions that vendor benchmark tables do not always state fully. The YOLO lineage's persistent advantage in this literature is ecosystem maturity: mature export paths to ONNX, TensorRT, CoreML and TFLite \citep{ultralytics2024yolo11,jocher2026yolo26}, whereas DETR-family models' deformable-attention operators have historically complicated embedded deployment. YOLO26's removal of both NMS and distribution focal loss is explicitly motivated by this deployment friction \citep{jocher2026yolo26}; whether the resulting CPU-speed claims hold on the specific hardware profile of a municipal deployment remains to be measured rather than assumed.

\section{Promptable Segmentation: The Segment Anything Line}
\label{sec:lr-sam}

The Segment Anything Model (SAM) introduced promptable segmentation as a foundation-model task: an image encoder, a prompt encoder for points and boxes, and a fast mask decoder, trained on the SA-1B corpus of one billion masks \citep{kirillov2023segment}. SAM segments what a prompt indicates but does not name it, and its text-prompting capability remained exploratory. SAM~2 extended the design to video with a streaming memory mechanism and improved image segmentation speed \citep{ravi2024sam2}. SAM~3 changed the task itself: \emph{promptable concept segmentation} accepts a short noun phrase or image exemplar and returns boxes, masks and identities for \emph{all} matching instances in an image or video, implemented as an 848M-parameter model with a shared backbone, a DETR-style detection head, a decoupled presence head separating recognition from localisation, and a memory-based tracker, trained via a data engine yielding four million concept labels with hard negatives and reported to roughly double concept-segmentation accuracy relative to prior systems on the new SA-Co benchmark \citep{meta2025sam3}. SAM~3.1 is a drop-in checkpoint update of the same architecture, adding object multiplexing that tracks up to sixteen objects in one forward pass and doubling video throughput \citep{meta2026sam31}.

For municipal use, the SAM line's appeal is obvious---open-vocabulary text prompting with genuine per-box confidence scores and masks, with no annotation or training cost---but its published evaluation profile invites caution. SA-Co performance measures agreement with a web-scale concept distribution, not with a specific operational taxonomy: a prompt such as ``orange garbage bag'' must traverse the model's own concept space, and nothing guarantees alignment with the dataset's operational class definition (Section~\ref{sec:bg-mdwd}). Known limitations relevant here include small and low-contrast objects, prompt ambiguity, domain shift from web imagery to specialised scenes, cluttered instance separation, and the computational cost of a ~3.5\,GB model relative to any supervised detector in this study. Mask outputs also require conversion to boxes for comparison against box-annotated ground truth, which introduces its own error term.

\section{Grounded and Open-Vocabulary Detection}
\label{sec:lr-grounding}

A complementary line produces boxes directly from text. GLIP unified detection and phrase grounding by aligning region and word embeddings under joint pretraining \citep{li2022grounded}; OWL-ViT transferred contrastively pretrained vision-language backbones to open-vocabulary detection with lightweight heads \citep{minderer2022simple}, scaled substantially by self-training in OWLv2 \citep{minderer2023scaling}; and Grounding DINO married a DINO-style DETR detector with grounded pretraining, becoming the de facto open-set detector and, in combination with SAM (``Grounded SAM''), a standard pipeline for text-to-mask extraction \citep{liu2024grounding,ren2024grounded}. These systems are reviewed here as architectural context: they define the design space of text-conditioned localisation---separate grounding-plus-segmentation pipelines versus unified single-pass models such as SAM~3---but none of them is among the models evaluated in this dissertation. The distinction the terminology must preserve is that \emph{open-vocabulary detection} names categories, \emph{phrase grounding} localises arbitrary referring text, and \emph{promptable segmentation} produces masks from spatial or textual prompts; systems differ in which of these they perform natively and in whether their confidence scores are comparable across prompts, a recurring evaluation difficulty in this literature.

\section{LocateAnything}
\label{sec:lr-locateanything}

LocateAnything-3B (NVIDIA, released mid-2026) is a dedicated visual-grounding vision-language model rather than a repurposed general VLM: a native-resolution MoonViT vision encoder feeds a Qwen2.5-3B language decoder, and localisation output is reformulated as fixed-length box-aligned blocks decoded in parallel---Parallel Box Decoding---rather than as a flat autoregressive coordinate stream, which the authors report yields up to 2.5$\times$ higher grounding throughput alongside state-of-the-art results across detection, phrase grounding, text localisation and pointing benchmarks \citep{nvidia2026locateanything}. The model is documented in an NVIDIA technical report and open-weight release rather than a peer-reviewed venue, so its benchmark claims carry vendor-report status. Compared with the SAM line it produces boxes and points but no masks; compared with grounded detectors such as Grounding DINO it inherits a generative VLM's flexibility in prompt phrasing but also its evaluation weakness: emitted boxes carry no calibrated per-box confidence, so ranking-based metrics such as AP are not meaningfully computable, and thresholded operating-point metrics must be used instead. Coordinates are emitted on a normalised grid following Qwen-VL conventions \citep{bai2025qwen25vl}. No published evaluation of LocateAnything on waste, signage or comparable municipal categories was identified.

\section{World Foundation Models and Cosmos Reason2}
\label{sec:lr-cosmos}

World models---systems that learn predictive models of an environment's dynamics---originate in reinforcement learning research \citep{ha2018world} and have been industrialised as \emph{world foundation models} (WFMs): large pretrained models of physical scenes intended to support simulation, prediction and embodied decision-making, of which NVIDIA's Cosmos platform is the most prominent open ecosystem \citep{nvidia2025cosmosplatform}. Within that ecosystem, the Cosmos Reason line provides the reasoning interface: Cosmos-Reason1 established a physically grounded VLM trained with supervised fine-tuning and reinforcement learning over physical-common-sense and embodied-reasoning ontologies \citep{nvidia2025cosmosreason1}, and Cosmos-Reason2 (released December 2025) scales this recipe onto the Qwen3-VL architecture in 2B, 8B and 32B variants, adding improved spatial--temporal understanding and---directly relevant here---structured visual grounding, in which the model emits bounding-box coordinates as JSON within its generated text \citep{nvidia2025cosmosreason2}. The three variants share architecture and differ in capacity and cost: roughly 5, 16 and 64\,GB of weights respectively at half precision, placing only the two smaller variants within single-consumer-GPU reach. Documentation-level claims that scaling improves reasoning quality are consistent with general VLM scaling evidence, but no published study isolates the effect of Cosmos-Reason2 scale on \emph{localisation} accuracy specifically.

Cosmos-Reason2 must not be described as a conventional detector. Its boxes are generated tokens produced through chain-of-thought reasoning, without a detection head, dense proposals or per-box confidence; output format compliance is prompt-dependent, and localisation competes with fluency objectives inside a generative decoder. Methodologically, this creates a genuine comparison problem that the evaluation literature has not resolved: supervised detectors are characterised by confidence-ranked precision--recall behaviour, whereas reasoning VLMs deliver an unranked set of assertions, so any cross-paradigm comparison must either collapse detectors to single operating points or forgo ranking metrics for the VLMs, and calibration analysis \citep{guo2017calibration} is unavailable for the latter. The attraction that justifies evaluating such models regardless is their zero-annotation adaptability and their capacity, in principle, to exploit contextual reasoning (``a bag placed at a kerb on collection day'') that no fixed-vocabulary detector represents.

\section{Self-Supervised Representation Learning and DINOv3}
\label{sec:lr-dinov3}

Self-supervised pretraining has progressed from contrastive instance discrimination \citep{chen2020simple,he2020momentum} through negative-free self-distillation \citep{grill2020bootstrap} and masked-image modelling \citep{he2022masked} to the DINO line, whose self-distillation over vision transformers produces features with notable emergent segmentation structure \citep{caron2021emerging}. DINOv2 scaled the recipe with curated data to features that rival supervised pretraining across classification, retrieval and dense tasks without fine-tuning \citep{oquab2024dinov2}. DINOv3 \citep{simeoni2025dinov3} scaled further---a 7B-parameter ViT trained on 1.7 billion curated images---and introduced Gram anchoring, a regulariser that constrains the pairwise-similarity structure of patch features to that of an earlier checkpoint, specifically to prevent the degradation of dense features over long training schedules; the released family distils this flagship into smaller ViT and ConvNeXt variants, including the ViT-B/16 scale used in this dissertation. The reported results show frozen DINOv3 features outperforming specialised state-of-the-art methods on dense benchmarks, which is the property most relevant to linear probing.

The critical question for this work is whether generic representation strength transfers to fine-grained, domain-specific attributes. Sign-condition assessment depends on subtle surface cues---rust texture, fading, sticker residue---that differ from the object-level semantics dominating SSL evaluation suites, and DINOv3's own evaluations do not cover deterioration-style attributes. Frozen-feature performance on such attributes is therefore an open question rather than a foregone conclusion, and it is exactly the question the dissertation's frozen-probe conditions operationalise. It bears repeating that DINOv3 the \emph{representation model} is unrelated to DINO the \emph{detector} \citep{zhang2023dino}, and that the dissertation uses only the former.

\section{Joint-Embedding Predictive Architectures and V-JEPA 2.1}
\label{sec:lr-vjepa}

The JEPA programme argues that predicting \emph{representations} of missing content, rather than reconstructing pixels, yields more semantic features \citep{lecun2022path}. I-JEPA instantiated this for images, predicting latent representations of masked regions from context \citep{assran2023self}; V-JEPA extended the objective to video by predicting masked spatiotemporal regions in feature space \citep{bardes2024revisiting}; V-JEPA~2 scaled to internet video, demonstrating strong motion understanding and downstream planning applications \citep{assran2025vjepa2}. The version used in this dissertation, V-JEPA~2.1 (March 2026), targets the family's acknowledged weakness---coarse, temporally unstable dense features---by supervising both masked and visible tokens with a distance-weighted dense predictive loss applied hierarchically across encoder layers, and unifies image and video training through multi-modal tokenisers \citep{murlabadia2026vjepa21}.

The relevance of video-pretrained features to \emph{static} sign-crop classification is the least-supported link in this dissertation's representation-learning chain, and the literature should be read as offering hypotheses rather than guarantees. V-JEPA-family evaluations concentrate on video benchmarks (action recognition, anticipation, temporal reasoning), with image-level results reported mainly to demonstrate breadth; published evidence that motion-derived features encode the fine surface statistics needed for condition assessment does not exist. Conversely, the objective's emphasis on predictable structure could plausibly favour geometry-related attributes such as shape and viewing angle. Applying a video encoder to still crops also carries mechanical costs---inputs must be presented as short clips, raising compute per crop---which is an operational disadvantage independent of representation quality.

\section{ConvNeXt and Modern Convolutional Networks}
\label{sec:lr-convnext}

ConvNeXt provides the supervised convolutional reference point. \citet{liu2022convnet} showed that a standard ResNet \citep{he2016deep}, modernised step by step with design choices associated with vision transformers---patchified stems, large depthwise kernels, inverted bottlenecks, LayerNorm, fewer activations---matches or exceeds comparable ViT accuracy on ImageNet and downstream tasks while remaining a pure convolutional network; ConvNeXt~V2 later added masked-autoencoder pretraining to the family \citep{woo2023convnextv2}. It is worth stating precisely that ConvNeXt is \emph{not} a CNN--Transformer hybrid: it contains no attention operations, and its significance lies in demonstrating that much of the ViT era's progress derived from training recipes and macro-design rather than attention itself. Compared with ViTs, ConvNeXt retains convolutional locality and hierarchical inductive biases, which improve sample efficiency on modest datasets and yield predictable, kernel-friendly computation on edge hardware; the ConvNeXt-Tiny scale used in this dissertation (~28M parameters, ImageNet-1k supervised pretraining) is thus both a widely deployed practical baseline and a deliberately different point in the pretraining space---supervised, image-based, convolutional---against the two self-supervised transformers reviewed above.

\section{LingBot-Vision and Boundary-Centric Dense Pretraining}
\label{sec:lr-lingbot}

LingBot-Vision (Robbyant/Ant Group, 2026) is a family of self-supervised vision transformers pretrained explicitly for \emph{dense spatial perception} rather than for general semantic recognition \citep{fu2026lingbot}. Its objective, \emph{masked boundary modelling}, is boundary-centric: a teacher--student scheme combines semantic self-distillation with a categorical boundary-field supervision signal, so that the learned patch features are optimised to preserve object boundaries and shapes. The released family spans ViT-S/16 to ViT-g/16, trained from random initialisation with rotary (RoPE) position embeddings and register tokens at a native pretraining crop size of 512~px; the authors recommend the normalised frozen patch tokens for dense readouts, and report backbone-level strength on boundary-sensitive downstream tasks such as depth estimation, semantic segmentation and video object segmentation. As with the newest detector generations reviewed above, the evidence base is a vendor technical report and official model release rather than an independent peer-reviewed evaluation, and no published study applies LingBot-Vision to traffic-sign or municipal-monitoring imagery.

Its inclusion in this dissertation's attribute study is a deliberate representational contrast rather than a claim of superiority. The three families reviewed so far span supervised classification features (ConvNeXt), general-purpose semantic self-distillation (DINOv3) and temporal predictive embedding (V-JEPA~2.1); none of their objectives is explicitly spatial or boundary-oriented. Two of the four MTSD attributes, however, are plausibly boundary-governed: sign \emph{shape} is literally an outline property, and \emph{physical condition} partly manifests as boundary disruption (bent edges, missing sections, holes). A boundary-centric representation therefore constitutes a directed hypothesis---that dense, shape-preserving pretraining transfers to outline- and damage-related attributes---which the frozen-probe and LoRA comparisons of the attribute study can test against the other three objectives under identical training conditions. Whether boundary-oriented features also help the context-dependent attributes (mounting, viewing angle), which depend on evidence \emph{around} the sign rather than its outline, is left open by the literature.

\section{Linear Probing, LoRA and Full Fine-Tuning}
\label{sec:lr-adaptation}

How much of a pretrained backbone to adapt is a research question in its own right. \emph{Linear probing} freezes the backbone and trains a shallow head; it is the standard instrument for measuring representation quality \citep{chen2020simple,oquab2024dinov2}, is extremely cheap in compute and memory, and cannot overwrite pretrained knowledge---but it also cannot correct any mismatch between pretrained features and the target domain, so it lower-bounds achievable accuracy under domain shift. \emph{Full fine-tuning} updates every parameter, maximising adaptation capacity, but on small datasets it risks overfitting and catastrophic forgetting of pretrained structure, and its memory footprint scales with backbone size \citep{yosinski2014transferable,he2019rethinking}. Between the extremes, parameter-efficient fine-tuning inserts small trainable modules into a frozen backbone \citep{houlsby2019parameter}; LoRA, the dominant such method, adds low-rank update matrices to selected weight projections---classically the attention query and value projections---training typically well under one per cent of parameters while approaching full fine-tuning quality on language tasks \citep{hu2022lora}. Its transfer to vision backbones is more recent and less settled: surveys of parameter-efficient adaptation note that most LoRA evidence derives from language and multimodal models, with vision-specific questions---rank selection, which projections to adapt, interaction with convolutional versus attention architectures---resolved only empirically per case \citep{han2024parameter}. Three consequences follow for experimental design. First, LoRA should not be presumed superior; its advantage over probing depends on how far the target domain lies from the pretraining distribution. Second, comparisons across adaptation strategies confound trainable capacity with strategy unless capacity is reported alongside results. Third, full fine-tuning of very large backbones may be infeasible on available hardware, which makes the probe-versus-LoRA comparison the operationally decisive one for foundation-scale encoders---precisely the comparison the dissertation's attribute study performs across its four backbone families and two model sizes.

\section{Multi-Task and Multi-Attribute Learning}
\label{sec:lr-multitask}

Hard parameter sharing---one backbone, several task heads---is the oldest and most deployment-friendly multi-task architecture \citep{caruana1997multitask,ruder2017overview}, with soft-sharing alternatives such as cross-stitch networks exchanging features between task-specific columns at higher cost \citep{misra2016cross}. The central empirical findings of the field are that sharing helps when tasks are related and data are limited, that \emph{negative transfer} between poorly matched tasks is real and can be severe \citep{standley2020tasks}, and that performance is sensitive to loss balancing, for which uncertainty weighting \citep{kendall2018multi}, gradient normalisation \citep{chen2018gradnorm} and gradient surgery \citep{yu2020gradient} are the standard remedies; dense-prediction surveys confirm that no single balancing scheme dominates \citep{vandenhende2021multi}. The closest applied literature is pedestrian-attribute recognition, which routinely predicts dozens of attributes from a shared representation and documents the same imbalance and correlation phenomena at scale \citep{wang2022pedestrian}. Traffic-sign-specific attribute learning is strikingly thin: the datasets reviewed in Section~\ref{sec:lr-signs} annotate identity rather than attribute bundles, and no published work identified in this search jointly predicts condition, mounting, viewing angle and shape for signs. The MTSD attribute setting is therefore near-novel territory, and the multi-task literature supplies design guidance (shared backbone, per-head class weighting, attention to imbalanced heads) rather than directly comparable baselines.

\section{Operational Comparison of Paradigms}
\label{sec:lr-operational}

The three families reviewed above occupy systematically different positions on the operational dimensions that determine municipal viability. Supervised detectors demand the full annotation pipeline---collection, labelling, quality assurance---and per-domain training, but reward this investment with compact models, millisecond latency, mature export tooling and calibrated confidence \citep{jocher2026yolo26,murshed2021machine}. Prompt-based foundation models invert the profile: near-zero annotation and training cost and unlimited class extensibility, against multi-gigabyte checkpoints, GPU-bound latency, prompt sensitivity, and---for the generative variants---absent confidence calibration \citep{meta2025sam3,nvidia2025cosmosreason2}. Representation-learning classifiers with light adaptation occupy a middle position, converting modest labelled sets into competitive task models at low training cost \citep{hu2022lora,oquab2024dinov2}. What the literature conspicuously lacks is joint measurement: studies reporting accuracy alongside annotation effort, latency, memory and maintenance burden for the \emph{same} task and data are rare, reproducibility of prompt-based results across prompt phrasings is seldom quantified, and label quality itself---which directly bounds every reported metric---is rarely audited, despite demonstrations that benchmark label errors materially distort model rankings \citep{northcutt2021pervasive} and standing calls for systematic dataset documentation \citep{gebru2021datasheets}. These omissions define the comparative and operational design of this dissertation's evaluation.

\section{Synthesis and Research Gap}
\label{sec:lr-gap}

Six gaps emerge consistently from this review. First, \emph{no Malta-specific evidence exists}: no identified study evaluates modern computer vision on Maltese street imagery for either waste or signage, despite the locally specific taxonomies, signage designs and environmental conditions documented in Chapter~\ref{ch:background}. Second, \emph{collection-stream waste detection is unserved}: existing waste datasets label materials or generic litter, not the bag-colour collection streams on which municipal operations actually run, and their geographic policy-dependence makes foreign datasets structurally unable to fill this role. Third, \emph{attribute-rich sign monitoring lacks datasets and baselines}: no established benchmark jointly annotates condition, mounting, viewing angle and shape, and traffic-sign attribute classification correspondingly lacks published baselines. Fourth, \emph{controlled cross-generation and cross-family detector evidence is limited}: YOLO11, YOLO12 and YOLO26 have not been compared under identical training protocols on any municipal domain in the published literature, vendor benchmarks dominate the evidence base for the newest generations, and comparisons between YOLO-family and DETR-family models rarely control for scale and pretraining. Fifth, \emph{zero-shot foundation models are unevaluated against fixed municipal ground truth}: the SAM~3 line, LocateAnything and Cosmos-Reason2 have no published evaluation on operational taxonomies of this kind, and the methodological problem of scoring unranked generative localisation against confidence-ranked detection remains largely undiscussed. Sixth, \emph{operational and integrity dimensions are under-reported}: annotation effort, latency and memory on stated hardware, prompt sensitivity, annotation-quality auditing and split-leakage analysis are seldom reported together with accuracy, though each conditions real-world adoption.

These gaps map directly onto the dissertation's research questions: the controlled multi-generation, multi-family supervised benchmark on MDWD (RQ1); the fixed-protocol evaluation of promptable and reasoning models against supervised ground truth on both datasets (RQ2); the adaptation-strategy comparison across pretrained representation families for multi-attribute sign classification (RQ3); the joint measurement of operational trade-offs across paradigms (RQ4); and the explicit treatment of dataset integrity, annotation quality and statistical uncertainty (RQ5). The reviewed literature provides the architectural components and the methodological warnings for each of these; what it does not provide---and what the following chapters construct---is a directly comparable, integrity-audited benchmark of all three paradigms on the two Maltese municipal domains.
